\textbf{Background.} Diabetic retinopathy (DR) is graded from lesions, but grading, lesion segmentation and lesion detection are usually solved by separate models, and grade predictions are not required to agree with lesion evidence or accompanied by guarantees on how often they are right.
\textbf{Objective.} We study whether grounding the DR grade in zone-wise lesion evidence from a segmentation network, together with ordinal soft targets, soft-logic clinical-rule consistency and conformal referral sets, improves grading, and we report lesion segmentation and detection on the same data.
\textbf{Methods.} On the DDR dataset (\num{13673} fundus images, 757 with lesion masks and boxes), a U-Net produces four lesion maps that a differentiable bottleneck summarises into 24 evidence tokens; a small transformer combines them with a frozen image embedding to predict cumulative ordinal probabilities. All experiments ran on a four-core CPU. A leak-controlled grading test set (3633 images) was used because the lesion and grading partitions overlap.
\textbf{Results.} Lesion evidence raised the quadratic-weighted kappa from \nBoneQWK--\nBtwoQWK\ (embedding only) to \nMainQWK\ (95\% CI \nMainCIlo--\nMainCIhi), with AUC for referable DR \nMainAUC. This exceeded a fine-tuned ResNet-18 reference (\nFTQWK). Soft targets lowered calibration error to \nMainECE\ but the rule-consistency loss gave no measurable gain. Segmentation reached a mean AUPR of \nSegAUPRFtlbg\ (microaneurysms $\le0.17$) and detection a mean AP$_{0.3}$ of \nDetMapFtlbg. Conformal sets reached nominal coverage when calibration and test data were exchangeable (\nExchCovTen\ at 0.90) but not across the official DDR validation and test partitions (\nConfCovTen), which differ in difficulty. Zero-shot transfer to a $224\times224$ external set failed (kappa $\le0.23$).
\textbf{Conclusions.} Lesion-derived evidence is a clear, cheap source of grading information; the other components are modest, and generalisation and partition shift remain open problems.
